% !TEX program = pdflatex
% Cheat Sheet FORMULE (mini) -- Fondamenti di Elettronica 085746
% Ultra-sintetico: solo formule + micro-cue. Formule dalle schede tde-percorsi-mentali
% (+ box base da CSFdE.tex). Stile-box CSFdE (landscape 3 col), box NON spezzabili,
% color coding per famiglia (tinte vicine = stesso esercizio d'esame).
\documentclass[11pt]{extarticle}
\usepackage[a4paper,landscape,margin=0.7cm,top=0.7cm,bottom=0.6cm]{geometry}
\usepackage[utf8]{inputenc}
\usepackage[T1]{fontenc}
\usepackage[italian]{babel}
\usepackage{amsmath,amssymb,mathtools}
\usepackage{multicol}
\usepackage{xcolor}
\usepackage[most]{tcolorbox}
\usepackage{microtype}

\definecolor{CustomYellow}{RGB}{233,215,0}
\definecolor{cue}{gray}{0.42}

% ---- color coding per famiglia (tinte vicine = stesso esercizio) ----
\definecolor{cDiodi}{RGB}{46,139,87}   % ES1  verde
\definecolor{cRC}{RGB}{0,140,150}      % ES1  teal
\definecolor{cMos}{RGB}{40,95,200}     % ES2  blu
\definecolor{cLog}{RGB}{112,72,190}    % ES2  indaco
\definecolor{cOp}{RGB}{224,122,20}     % ES3  arancio
\definecolor{cAdc}{RGB}{200,52,72}     % ES3  rosso

\newcommand{\vov}{V_{\mathit{OV}}}

% ---- box formula: look CSFdE, NON breakable (sta in colonna), colore per famiglia ----
\newtcolorbox{fmbox}[2][CustomYellow]{%
  enhanced, colback=white, colframe=#1,
  boxrule=0.9pt, arc=2.5pt, left=3pt, right=3pt, bottom=3pt, top=6pt,
  title={#2}, coltitle=white, colbacktitle=#1,
  fonttitle=\bfseries\footnotesize,
  attach boxed title to top left={xshift=6pt, yshift=-3pt},
  boxed title style={colframe=#1, arc=1.5pt, boxrule=0pt},
  before skip=2pt, after skip=6pt}

\newcommand{\F}[2]{\(\displaystyle #1\)\;{\footnotesize\color{cue}#2}\par\vspace{3pt}}
\newcommand{\SH}[1]{\vspace{3pt}{\bfseries\small #1}\par\vspace{1.5pt}}
\newcommand{\SY}[1]{\vspace{2pt}\hrule height 0.2pt\vspace{2pt}{\footnotesize\color{cue}\textbf{Simboli:} #1}\par}

\setlength{\parindent}{0pt}
\setlength{\columnsep}{11pt}
\setlength{\columnseprule}{0pt}
\setlength{\parskip}{2pt}

\begin{document}
\begin{center}
  {\Large\bfseries Cheat Sheet Formule}\;\;{\large Fondamenti di Elettronica 085746}\\[1pt]
  {\footnotesize conv. Acconcia $K{=}\tfrac12\mu C'_{ox}\tfrac{W}{L}$ \;\textbullet\; formule + contesto essenziale \;\textbullet\; colore ${=}$ famiglia: {\color{cDiodi}\textbf{diodi}}/{\color{cRC}\textbf{RC}} (ES1) \, {\color{cMos}\textbf{MOS}}/{\color{cLog}\textbf{logica}} (ES2) \, {\color{cOp}\textbf{opamp}}/{\color{cAdc}\textbf{conv.}} (ES3) \;\textbullet\; \today}
\end{center}
\vspace{2pt}

\begin{multicols*}{3}

% ================= BASE =================
\begin{fmbox}{BASE — reti, dB, prefissi}
\SH{Reti (usate ovunque)}
\F{V_{R_x}=V\,\tfrac{R_x}{\sum R_i}}{partitore di tensione}
\F{I_{R_1}=I\,\tfrac{R_2}{R_1+R_2}}{partitore di corrente}
\F{R_s=\textstyle\sum R_i;\quad R_1\!\parallel\!R_2=\tfrac{R_1R_2}{R_1+R_2}}{serie / parallelo}
\F{\text{gen.\ spenti: }V\!\to\!\text{corto},\ I\!\to\!\text{aperto}}{sovrapposizione}
\F{V_{th}{=}V\text{ a vuoto},\ R_{th}{=}R\text{ vista}}{Thevenin (gen.\ spenti)}
\SH{dB \& prefissi SI}
\F{|G|_{dB}=20\log_{10}|G|}{$6$ dB$\,{\approx}\,2\times$, $20$ dB$\,{=}\,10\times$, $40$ dB$\,{=}\,100\times$}
\F{\text{m}\,10^{-3}\ \ \mu\,10^{-6}\ \ \text{n}\,10^{-9}\ \ \text{p}\,10^{-12}}{k M G: $10^{3},10^{6},10^{9}$}
\end{fmbox}

% ================= DIODI 1 =================
\begin{fmbox}[cDiodi]{DIODI — modelli \& metodo}
\SH{Modelli PWL ($V_\gamma\!\approx\!0{,}7$ V)}
\F{V_D=V_A-V_K,\ I_D>0\text{ diretta}}{convenzione segni}
\F{\text{ON}:\,V_D=V_\gamma,\ \text{ver. }I_D>0}{soglia diretta}
\F{\text{OFF}:\,I_D=0,\ \text{ver. }V_D<V_\gamma}{interdizione}
\F{\text{BD}:\,V_D=-|V_{BD}|,\ \text{ver. }I_D<0}{breakdown / zener}
\F{\text{zener OFF}:\,-V_Z<V_D<V_\gamma}{range interdizione}
\SH{Metodo a ipotesi}
\F{\text{ON}\!\to\!{-}V_\gamma,\ \text{OFF}\!\to\!\text{apre},\ \text{BD}\!\to\!{-}|V_{BD}|}{sostituisci, KVL/KCL}
\F{I_D=\tfrac{V_{in}-V_\gamma}{R}}{diodo ON in serie a $R$}
\F{V_D{=}V_\gamma\ \text{o}\ V_D{=}{-}|V_{BD}|}{$V_{in}$ di transizione}
\end{fmbox}

% ================= DIODI 2 =================
\begin{fmbox}[cDiodi]{DIODI — caratteristica \& potenza}
\SH{Clipper a massa}
\F{|V_{in}|<V_\gamma:\ V_{out}=V_{in}}{lineare (2 OFF)}
\F{V_{in}>{+}V_\gamma:\ V_{out}={+}V_\gamma}{clamp alto}
\F{V_{in}<{-}V_\gamma:\ V_{out}={-}V_\gamma}{clamp basso}
\SH{Potenza}
\F{P_D^{med}=\tfrac1T\!\int_0^T\! V_D I_D\,dt}{media su periodo}
\F{P_D^{med}=d\,V_\gamma I_{D,ON}}{onda quadra, duty $d$}
\SH{Raddrizz. SS $+C$ / rilev. picco / zener}
\F{\Delta V\approx V_p\tfrac{T}{R_LC}}{ripple pp ($R_LC\!\gg\!T$)}
\F{V_{BD,min}\gtrsim 2V_p,\quad V_{out}\!\approx\!V_p}{inversa / picco}
\F{I_{D,med}\approx\tfrac{V_{out}}{R_L},\quad I_{D,pk}=C\tfrac{dV}{dt}\big|_{cond}}{correnti diodo}
\F{R_S^{max}=\tfrac{V_{in,min}-V_Z}{I_Z^{min}+I_{L,max}}}{zener: worst-case}
\SY{$V_\gamma$ soglia \textbullet\ $|V_{BD}|{=}V_Z$ \textbullet\ $R_S$ zavorra \textbullet\ $V_p$ picco \textbullet\ $d$ duty}
\end{fmbox}

% ================= TRANSITORI RC =================
\begin{fmbox}[cRC]{TRANSITORI RC (1$^\circ$ ordine)}
\SH{Formula master}
\F{V_C(t)=V_\infty+(V_0-V_\infty)e^{-t/\tau}}{legge universale}
\F{\tau=R_{eq}C}{$R_{eq}$ vista dai morsetti di $C$}
\F{V_C(0^+)=V_C(0^-)}{continuità $\Rightarrow V_0$}
\SH{Ricavare $\tau$ e $t$}
\F{H(s)=\tfrac{K}{1+s\tau},\ \tau=\tfrac{1}{2\pi f_p}}{$\tau$ = coeff.\ $s$ al denom.}
\F{t=-\tau\ln\tfrac{V-V_\infty}{V_0-V_\infty}}{inverti la master}
\F{t\approx0{,}69\,n\,\tau}{soglia $<1$ LSB, ADC $n$ bit}
\F{V_C(t)=V_0+\tfrac{I}{C}\,t}{a corrente costante}
\SH{$C$ alla transizione / costanti}
\F{V_C{=}V_0\!\Rightarrow\text{gen. }V_0;\ t\!\to\!\infty\!\Rightarrow\!\text{aperto}}{corto se $V_0{=}0$}
\F{\text{gradino: }s\!\to\!\infty\Rightarrow\tfrac{1}{sC}\!\to\!0}{$C$ = corto per il $\Delta$}
\F{V(t_0^+){=}V(t_0^-){+}\Delta V_{in}\!\cdot\!\text{part.}|_{C\,\text{corto}}}{2° gradino: solo il $\Delta$, poi rilassa}
\F{0{,}69\tau\!:\!50\%,\ 1\tau\!:\!63\%,\ 2{,}2\tau\!:\!10\text{–}90\%}{$5\tau$: "stessa forma"}
\SY{$V_0$ iniziale \textbullet\ $V_\infty$ regime \textbullet\ $\tau{=}R_{eq}C$}
\end{fmbox}

% ================= MOSFET 1 =================
\begin{fmbox}[cMos]{MOSFET — modello \& zone}
\SH{Notazione unificata ($V^\star$ già $>0$)}
\F{V_{GS}^\star{=}V_{GS}\text{(n)},\,V_{SG}\text{(p)}}{gate-source utile}
\F{V_{DS}^\star{=}V_{DS}\text{(n)},\,V_{SD}\text{(p)}}{drain-source utile}
\F{\vov=V_{GS}^\star-|V_T|}{overdrive (da $V_{GS}^\star$!)}
\F{K=\tfrac12\mu C'_{ox}\tfrac{W}{L}}{conv. Acconcia}
\SH{ON/OFF (gate) \& zona (drain)}
\F{V_{GS}^\star>|V_T|\Rightarrow\text{ON}}{altrimenti $|I_D|{=}0$}
\F{V_{DS}^\star<\vov\Rightarrow\text{ohmica}}{sotto pinch-off}
\F{V_{DS}^\star>\vov\Rightarrow\text{saturazione}}{pinch-off}
\F{V_{DS}^\star=\vov\ \Leftrightarrow\ V_{GD}^\star{=}|V_T|}{confine ohmica/sat}
\end{fmbox}

% ================= MOSFET 2 =================
\begin{fmbox}[cMos]{MOSFET — corrente, $R$, S/D, PTL}
\SH{Corrente (universale col modulo)}
\F{|I_D|=K[2\vov V_{DS}^\star-(V_{DS}^\star)^2]}{ohmica}
\F{|I_D|=K\vov^2}{saturazione}
\F{I_{DS}={+}|I_D|\text{(n)},\,{-}|I_D|\text{(p)}}{segno finale (p: $S\!\to\!D$)}
\SH{$R$/gen. \& $K$ serie/parallelo}
\F{R_{DS,on}=\tfrac{1}{2K\vov}}{ohmica profonda $V_{DS}^\star\!\ll\!\vov$}
\F{I_{sat}=K\vov^2}{sat = gen.\ corrente}
\F{\tfrac{1}{K_{eq}}{=}\!\sum\tfrac{1}{K_i}\text{ (serie)},\ K_{eq}{=}\!\sum K_i\text{ (par.)}}{come conduttanze}
\SH{Source/Drain \& pass-transistor}
\F{\text{n: }S{=}V_{min};\ \text{p: }S{=}V_{max}}{$I$ da $V$ alto a basso}
\F{\text{n: }V_{C,max}{=}V_{DD}{-}V_T}{PTL: carica male gli alti}
\F{\text{p: }V_{C,min}{=}|V_T|}{PTL: scarica male i bassi}
\SY{$V_T{>}0$ (n), ${<}0$ (p) \textbullet\ TG${=}$n$\parallel$p rail-to-rail \textbullet\ Lusardi $K_L{=}2K$}
\end{fmbox}

% ================= PORTE 1 =================
\begin{fmbox}[cLog]{PORTE LOGICHE — stati \& cross-cond.}
\SH{4 stati d'uscita (PUN / PDN)}
\F{\text{PU on, PD off}\Rightarrow V_{OUT}{=}V_{DD}}{logico 1}
\F{\text{PU off, PD on}\Rightarrow V_{OUT}{=}0}{logico 0}
\F{\text{off/off}\Rightarrow\text{HZ},\ V_{OUT}{=}V_{OUT}(t^-)}{$C_L$ memorizza}
\F{\text{on/on}\Rightarrow\text{cross-conduzione}}{$V_{OUT}$ intermedia}
\SH{Cross-conduzione}
\F{K_{max}\!\to\!\text{ohm},\ K_{min}\!\to\!\text{sat}}{shortcut K-dominante}
\F{V_{OUT}=I_{sat}^{(K_{min})}R_{DS,on}^{(K_{max})}}{se $K_{max}/K_{min}\!\gg\!1$}
\F{K_pV_{OV\!,p}^2{=}K_n[2V_{OV\!,n}V_{DS,n}{-}V_{DS,n}^2]}{bilancio esatto (P sat, N ohm)}
\F{V_{OUT}{=}\tfrac{K_pV_{OV\!,p}^2}{2K_nV_{OV\!,n}}}{caso $K_n{>}K_p$ (logico 0)}
\F{V_{DD}{-}V_{OUT}{=}\tfrac{K_nV_{OV\!,n}^2}{2K_pV_{OV\!,p}}}{caso $K_p{>}K_n$ (logico 1)}
\end{fmbox}

% ================= PORTE 2 =================
\begin{fmbox}[cLog]{PORTE LOGICHE — tempi \& potenze}
\SH{Tempi propagazione ($V_{sog}{=}V_{DD}/2$)}
\F{t_p=\tfrac{C_L\Delta V}{I_{sat}},\ \Delta V{=}V_C(0){-}V_{sog}}{regime saturazione}
\F{t_p=\tau\ln\tfrac{V_C(0)-V_C(\infty)}{V_{sog}-V_C(\infty)}}{ohmico, $\tau{=}R_{DS,on}C$ ($+R$ est.)}
\F{R_{DS,on}{=}\tfrac{1}{2K\vov};\ \ R_{pinch}{=}\tfrac{\vov}{I_{sat}}}{tangente / corda ($=2R_{DS,on}$)}
\SH{PTL (valore finale ridotto)}
\F{\text{n (salita): }V_C(\infty)=V_{DD}{-}V_{Tn}}{carica male gli alti}
\F{\text{p (discesa): }V_C(\infty)=|V_{Tp}|}{scarica male i bassi}
\F{V_{sog}\text{ mai raggiunta}\Rightarrow t_p\!\to\!\infty}{n: $V_{sog}{>}V_C(\infty)$; p: $V_{sog}{<}V_C(\infty)$}
\SH{Potenze}
\F{P_{stat}=\tfrac{T_{cross}}{T}V_{DD}I_{cross}}{$0$ senza cross-cond.}
\F{P_{din}=N_{sal}\,f\,C\,\Delta V\,V_{DD}}{per nodo, poi somma}
\F{\Delta V:\ V_{DD}\ /\ V_{DD}{-}V_{Tn}\ /\ V_{DD}{-}|V_{Tp}|}{CMOS / PTLn / PTLp!}
\F{\textstyle\sum\Delta V_\uparrow=\sum\Delta V_\downarrow}{sanity check, se no errore}
\F{T_{patt}=\text{mcm dei periodi}}{segnali con $T$ diversi}
\SY{$V_{sog}{=}V_{DD}/2$ \textbullet\ $N_{sal}$ fronti salita/periodo \textbullet\ $I_{cross}{=}K_{min}\vov^2$}
\end{fmbox}

% ================= OPAMP 1 =================
\begin{fmbox}[cOp]{OPAMP — guadagni ideali \& $A(s)$}
\SH{$G_{id}$ (corto virt.\ $V^+{=}V^-$, $i_\pm{=}0$)}
\F{G_{id}=-\tfrac{R_2}{R_1}}{invertente}
\F{G_{id}=1+\tfrac{R_2}{R_1}}{non invertente}
\F{G_{id}=1}{buffer}
\F{G_{id}=-\tfrac{1}{sR_1C}}{integratore (polo orig.)}
\F{G_{id}=-sR_2C}{derivatore (zero orig.)}
\F{G_{id}=-\tfrac{R}{1+sRC}}{transimpedenza}
\SH{OpAmp reale (1 polo)}
\F{A(s)=\tfrac{A_0}{1+s\tau_0}}{anello aperto}
\F{f_0=\tfrac{1}{2\pi\tau_0},\quad \mathrm{GBWP}=A_0 f_0}{polo / dato di targa}
\F{|A|\approx\tfrac{\mathrm{GBWP}}{f}}{per $f\!\gg\!f_0$}
\end{fmbox}

% ================= OPAMP 2 =================
\begin{fmbox}[cOp]{OPAMP — impedenze, $G_{loop}$, $f^\ast$}
\SH{Impedenze \& forma di Bode}
\F{R\!\parallel\!\tfrac{1}{sC}=\tfrac{R}{1+sRC}}{parallelo RC (polo)}
\F{R+\tfrac{1}{sC}=\tfrac{1+sRC}{sC}}{serie RC (zero)}
\F{\text{cost}+s\,a=\text{cost}(1+s\tfrac{a}{\text{cost}})}{$\Rightarrow\tau{=}a/\text{cost}$}
\F{f_p=\tfrac{1}{2\pi\tau}}{polo/zero (occhio al $2\pi$)}
\SH{Conservazione guadagno$\times$banda (per tratto)}
\F{{-}20\text{: }|G|{\cdot}f{=}\text{cost};\quad {+}20\text{: }\tfrac{|G|}{f}{=}\text{cost}}{rapp.\ moduli $=$ rapp.\ freq.}
\F{{\pm}40\text{: }|G|{\cdot}f^{\pm2}{=}\text{cost}}{SOLO le $f$ al quadrato!}
\F{\text{tratto piatto: }|G|\text{ costante}}{l\`i non si taglia mai}
\SH{$G_{loop}$ (rompi anello, $V_{test}$ in uscita)}
\F{G_{loop}'=\tfrac{V_+-V_-}{V_{test}}}{fattore retroazione}
\F{G_{loop}=G_{loop}'\cdot A(s)}{anello completo}
\F{|G_{loop}(j2\pi f^\ast)|=1}{$f^\ast$: incrocio $|G_{id}|$--$|G_{OPEN}|$}
\F{f^\ast=|G_{loop}'(0)|\,\mathrm{GBWP}}{se poli/zeri $\gg f^\ast$}
\end{fmbox}

% ================= OPAMP 3 =================
\begin{fmbox}[cOp]{OPAMP — $G_{reale}$ \& stabilità}
\SH{$G_{reale}$ (metodo grafico)}
\F{|G_{reale}|=\min\{|G_{id}|,|G_{OPEN}|\}}{in dB; $G_{OPEN}{=}{-}G_{id}G_{loop}$}
\F{f\!\ll\!f^\ast\!:G_{id};\quad f\!\gg\!f^\ast\!:G_{OPEN}}{transizione a $f^\ast$}
\F{G_{reale}\approx\tfrac{G_{id}}{1+s\tau^\ast},\ \tau^\ast{=}\tfrac{1}{2\pi f^\ast}}{polo dominante}
\SH{Stabilità (su $G_{loop}$, tutti poli/zeri)}
\F{\phi_m=180^\circ{-}\!\sum_p\!\arctan\tfrac{f^\ast}{f_p}{+}\!\sum_z\!\arctan\tfrac{f^\ast}{f_z}}{margine di fase}
\F{\text{1 polo}{<}f^\ast\!:\ {\sim}90^\circ;\quad\text{polo a }f^\ast\!:\ 45^\circ}{regola pratica}
\F{\phi_m\gtrsim45^\circ\text{ stabile};\quad \phi_m\to0^\circ\text{ oscilla}}{2 poli vicini: instab.}
\F{\text{taglio a }{-}40\text{ dB/dec}\Rightarrow\text{instabile}}{scorciatoia TdE}
\SY{$s{=}j2\pi f$ \textbullet\ polo ${-}20$, zero ${+}20$ dB/dec \textbullet\ $G_{loop}'$ fattore reaz.}
\end{fmbox}

% ================= OPAMP non idealità =================
\begin{fmbox}[cOp]{OPAMP — non idealità (DC + SR)}
\SH{Procedura: ideale $+$ $C$ aperti, sovrapposizione}
\F{V_{out}=\pm(1+\tfrac{R_f}{R_g})V_{os}}{offset $\times$ guad.\ rumore}
\F{V_{out}=R_f I_{b-}}{bias, contributo pin $-$}
\F{V_{out}=(1+\tfrac{R_f}{R_g})R_+I_{b+}}{$R_+$ vista dal $+$}
\F{R_+{=}R_f\!\parallel\!R_g\Rightarrow V_{out}=R_f I_{os}}{con compensazione}
\F{I_{os}=|I_{b+}-I_{b-}|}{corrente di offset}
\F{V_{out}^{wc}=\textstyle\sum|\text{contributi}|}{segni ignoti: somma moduli}
\SH{Slew rate}
\F{2\pi f V_p\le \mathrm{SR}}{no distorsione (sinusoide)}
\F{\mathrm{FPBW}=\tfrac{\mathrm{SR}}{2\pi V_p}}{full-power bandwidth}
\SY{$V_{os}$ offset tens. \textbullet\ $I_{b\pm}$ bias \textbullet\ $\mathrm{SR}{=}\max|\dot V_{out}|$ \textbullet\ ingressi DC: SR ininfluente}
\end{fmbox}

% ================= ADC & DAC =================
\begin{fmbox}[cAdc]{ADC \& DAC}
\SH{ADC fondamentali}
\F{1\,\mathrm{LSB}=\tfrac{V_{FS}}{2^n}}{risoluzione / scalino}
\F{V_{FS}=\min(V_{FS}^{ADC},\Delta V_{out}^{op})}{fondo scala effettivo}
\F{\mathrm{LSB}_{in}=\tfrac{\mathrm{LSB}_{ADC}}{G_{id}}}{risoluzione all'ingresso}
\F{\varepsilon_q=\pm\tfrac12\mathrm{LSB}}{errore quantizzazione}
\SH{Tempi di conversione (cicli $\times T_{clk}$)}
\F{t_c^{SAR}=n\,T_{clk}}{($+1$ con S/H)}
\F{t_c^{grad/inseg}=2^n T_{clk},\quad t_c^{rampa}=2^n T_{clk}}{fino a $2^n$ cicli}
\F{t_c^{2\text{ rampe}}=2^{n+1}T_{clk}}{integra $+$ deintegra}
\F{2^n\ge N_{max}=T_{int}/T_{clk}}{bit contatore (doppia rampa)}
\F{f_{clk}^{min}{=}\{n,\,2^n,\,2\!\cdot\!2^n\}/T_{hld}}{SAR / rampa / doppia}
\SH{DAC}
\F{V_{out}={-}\tfrac{V_{ref}}{2}\big(\tfrac{D_i}{2^0}{+}\tfrac{D_{i-1}}{2^1}{+}\dots\big)}{R-2R, MSB peso max}
\F{\text{DNL, INL, guadagno, offset}}{monotono se DNL$\ge{-}1$ LSB}
\end{fmbox}

% ================= SAMPLE & HOLD =================
\begin{fmbox}[cAdc]{Sample \& Hold}
\SH{$\tau$ e vincoli}
\F{\tau_{smp}{=}R_{DS,on}C_H,\quad \tau_{hld}{=}R_{ADC}C_H}{$R$ vista da $C_H$}
\F{T_{smp}>0{,}69\,n\,\tau_{smp}}{acquisizione $<1$ LSB}
\F{T_{hold}\ge T_{conv}^{ADC}}{copre la conversione}
\F{V_G>\vov^{min}+V_t+V_{S,max}}{gate sample (worst $V_S$)}
\F{C_H\ge\tfrac{T_{hld}(2^n-1)}{R_{ADC}}}{dimensionam.\ (sfrutta dinamica)}
\SH{Errori}
\F{V_{inj}=\Delta V_G\tfrac{C_{par}}{C_{par}+C_H}}{charge injection ($<\tfrac12$ LSB)}
\F{\Delta V_{drop}=\tfrac{I_b}{C_H}T_{hld}}{droop/drop in hold}
\F{V_\varepsilon=V_{in}\tfrac{G_{id}}{1+|G_{loop}|}}{err.\ statico di guadagno (DC)}
\F{E_{fs}=V_{in,fs}(G_{id}-G_{reale})}{errore a fondo scala}
\SY{$V_{FS}$ fondo scala \textbullet\ $C_H$ hold \textbullet\ $C_{par}$ parassita \textbullet\ $I_b$ leakage}
\end{fmbox}

\end{multicols*}
\end{document}
